\section{Sub-CPMK / Indikator Pencapaian}

Sub-CPMK adalah penjabaran CPMK menjadi indikator yang lebih terukur dan dapat diuji. Setiap CPMK dijabarkan menjadi beberapa Sub-CPMK yang spesifik, sehingga asesmen dapat dirancang untuk mengukur pencapaian setiap indikator secara jelas. Dalam buku ini, setiap bab memetakan materinya ke Sub-CPMK yang relevan, sehingga Anda dapat melacak capaian belajar secara terstruktur \cite{mdnLearn}.

Penjabaran CPMK menjadi indikator yang lebih terukur dan dapat diuji:

\textbf{Turunan CPMK-1}
\begin{itemize}[leftmargin=*]
  \item \textbf{Sub-CPMK 1.1:} Mengidentifikasi kebutuhan fungsional dan non-fungsional pengguna untuk solusi web.
  \item \textbf{Sub-CPMK 1.2:} Membuat wireframe dan mockup antarmuka web yang mendukung alur pengguna (user flow).
  \item \textbf{Sub-CPMK 1.3:} Menerapkan prinsip aksesibilitas (a11y) dan pengalaman pengguna (UX) dalam desain antarmuka.
\end{itemize}

\textbf{Turunan CPMK-2}
\begin{itemize}[leftmargin=*]
  \item \textbf{Sub-CPMK 2.1:} Membuat halaman web dengan HTML5 (struktur, elemen semantik, form, media).
  \item \textbf{Sub-CPMK 2.2:} Menerapkan CSS3 untuk styling, layout (Flexbox/Grid), dan desain responsif.
  \item \textbf{Sub-CPMK 2.3:} Mengimplementasikan interaktivitas dengan JavaScript (DOM, event handling, validasi form).
  \item \textbf{Sub-CPMK 2.4:} Menggunakan library atau framework front-end untuk antarmuka dinamis.
\end{itemize}

\textbf{Turunan CPMK-3}
\begin{itemize}[leftmargin=*]
  \item \textbf{Sub-CPMK 3.1:} Mengimplementasikan logika server-side dengan PHP (pemrosesan form, struktur dasar, OOP).
  \item \textbf{Sub-CPMK 3.2:} Mengelola basis data dan mengimplementasikan operasi CRUD dengan MySQL serta koneksi dari PHP.
  \item \textbf{Sub-CPMK 3.3:} Mengintegrasikan API dan menghubungkan front-end dengan back-end (request/response, data terstruktur).
\end{itemize}

\textbf{Turunan CPMK-4}
\begin{itemize}[leftmargin=*]
  \item \textbf{Sub-CPMK 4.1:} Melakukan pengujian dan evaluasi performa aplikasi web.
  \item \textbf{Sub-CPMK 4.2:} Menerapkan prinsip keamanan web: validasi input, prepared statement, session, mitigasi XSS dan CSRF.
  \item \textbf{Sub-CPMK 4.3:} Melakukan deployment ke lingkungan produksi dan administrasi dasar (server/web, basis data).
\end{itemize}

Sub-CPMK di atas dipetakan ke 16 pertemuan dalam satu semester. Tabel berikut menyajikan rencana pembelajaran lengkap yang memetakan setiap pertemuan ke Sub-CPMK, materi pokok, bentuk pembelajaran, dan asesmen.
